\item Pierre de Fermat ($1601-1665$) demostró que la trayectoria de un rayo de luz entre dos puntos es aquella que requiere el intervalo de tiempo mínimo (Principio de Fermat). Deduzca la ley de refracción de Snell a partir del principio de Fermat usando la figura P34.54:
\begin{enumerate}
    \item Demuestre que el tiempo que tarda la luz en ir de P a Q es:
    $$ t = \frac{r_1}{v_1} + \frac{r_2}{v_2} = \frac{n_1 \sqrt{a^2 + x^2}}{c} + \frac{n_2 \sqrt{b^2 + (d - x)^2}}{c} $$
    \item Para obtener el valor de $x$ para el cual $t$ tiene su valor mínimo, derive $t$ respecto a $x$ e iguale la derivada a cero. Demuestre que el resultado implica:
    $$ \frac{n_1 x}{\sqrt{a^2 + x^2}} = \frac{n_2 (d - x)}{\sqrt{b^2 + (d - x)^2}} $$
    \item Demuestre que esta expresión da directamente la ley de Snell: $n_1 \sin\theta_1 = n_2 \sin\theta_2$.
\end{enumerate}